Units $\hbar = c = k_{B} = 1$ except where cm and s are written; $3 < \pi < 22/7$ wherever $\pi$ enters a bound; $M_{P} = 2.435 \times 10^{18}$ GeV is the reduced Planck mass.

\emph{Literature inputs: premises of the statement, not imported in the store, cited from memory (to be imported as Facts).} (a) For a simply connected group $G$ broken to $H$, $\pi_{2}(G/H) \cong \pi_{1}(H)$ (the exact homotopy sequence of $H \to G \to G/H$), and every element of $\pi_{2}(G/H)$ is carried by a finite-energy 't Hooft-Polyakov monopole; for Pati-Salam broken to the Standard Model the lightest monopole carries two Dirac units (Lazarides, Magg and Shafi 1980). (b) The Kibble mechanism with Preskill's estimate: a broken phase formed with an orientation uncorrelated beyond a length $\xi \le H^{-1}$ leaves at least about $p \approx \tfrac{1}{10}$ monopole per $\xi^{3}$ (Kibble 1976, Preskill 1979); (K) uses only $\xi \le H^{-1}$, which causality enforces, and $p \ge \tfrac{1}{10}$. (c) Preskill's annihilation: monopole-antimonopole annihilation in the plasma reduces $n_{M}/s$ only to a floor of about $10^{-10}(M_{M}/10^{14}\,\mathrm{GeV})$ and acts only above it. (d) The MACRO final flux limit for slow and fast monopoles, $\mathcal{F} < 1.4 \times 10^{-16}$ cm$^{-2}$ s$^{-1}$ sr$^{-1}$ over $4 \times 10^{-5} \lesssim \beta < 1$ (MACRO Collaboration 2002); galactic monopoles move at least at the halo's virial speed, $\beta \ge 7 \times 10^{-4}$, and light ones are accelerated by the galactic field to larger $\beta$, where the IceCube limits are stronger still. (e) High-temperature symmetry restoration, with its known exceptions for negative cross-couplings (Weinberg 1974; Dvali, Melfo and Senjanovic 1995). Not a premise but a listed escape: the Langacker-Pi mechanism (Langacker and Pi 1980), in which a temporary breaking of electromagnetism at intermediate temperatures confines monopoles and antimonopoles in flux tubes and annihilates them; it needs an additional stage of the descended phase that the stored chain does not contain.

(i) In the series' language $G = SU(4) \times SU(2)_{L} \times SU(2)_{R} \cong \mathrm{Spin}(6) \times \mathrm{Spin}(4)$, the even Clifford group of the slot module that the rung parity leaves after the first breaking (Proposition \claimref{standard-model.majorana-element}), is simply connected. On the $4$ of $SU(4)$, $\tfrac12(B - L) = \operatorname{diag}(\tfrac16, \tfrac16, \tfrac16, -\tfrac12)$, traceless; $T_{3R} = \pm\tfrac12$ on the $SU(2)_{R}$ doublet, traceless; so $Y = T_{3R} + \tfrac12(B - L)$ lies in the semisimple Lie algebra of $G$ and $t \mapsto \exp(2\pi tY)$ (the phase unit realised as a unit bivector) is a closed one-parameter subgroup. On the $(4,2,1)$ it gives $\tfrac16$ and $-\tfrac12$, and on the conjugates of the $(\bar 4,1,2)$ it gives $-\tfrac23, \tfrac13, 0, 1$: the Standard Model hypercharges (Fact \claimref{standard-model.imported-pati-salam-representation-theory-survival}), with least common period $6$. It commutes with $SU(3) \times SU(2)_{L}$, so $H$ contains a central circle and $\pi_{1}(H) \supseteq \mathbb{Z}$; by (a) the breaking at $M_{I}$ leaves stable monopoles. Their mass is of order $4\pi v_{R}/g_{R}$, above $10^{12}$ GeV with $g_{R} = 0.469$ and $v_{R} = M_{I}/g_{R} \in (4.15, 4.16) \times 10^{10}$ GeV (Proposition \claimref{standard-model.what-angle-asks-neutrino-sector}).

(ii) In the radiation-dominated descended phase $H = \pi\sqrt{g_{*}/90}\,T^{2}/M_{P}$ and $s = (2\pi^{2}/45)g_{*}T^{3}$, so (K) gives $n_{M}/s \ge pH^{3}/s$, and $(pH^{3}/s)^{2} = p^{2}\pi^{6}(g_{*}/90)^{3}(T/M_{P})^{6}/((2\pi^{2}/45)^{2}g_{*}^{2}) = p^{2}\pi^{2}g_{*}(T/M_{P})^{6}/1440$ (the constant $45^{2}/(4 \cdot 90^{3}) = 1/1440$ is checked). The bound is increasing in $g_{*}$, so $g_{*} \ge 106.75$ makes it conservative. Under (D) there is no inflationary stage after the descent and entropy is conserved, so $n_{M}/s$ is today's value. Today $s_{0} = (2\pi^{2}/45)g_{s}(T_{0})T_{0}^{3}$ with $T_{0} = 2.7255$ K (Fact \claimref{synthesis.imported-planck-2018-lcdm-parameters}), $k_{B}T_{0}/\hbar c = 11.90$ cm$^{-1}$, $s_{0} = 2891$ cm$^{-3}$ (checked between $2600$ and $3200$ with the $\pi$ brackets). An isotropic flux $\mathcal{F}$ of monopoles of speed $\beta$ has $n_{M} = 4\pi\mathcal{F}/(\beta c)$, so (d) gives $n_{M}/s_{0} < 4\pi\mathcal{F}/(\beta c\,s_{0}) < 3.7 \times 10^{-26}$ at $\beta = 7 \times 10^{-4}$ (the cosmic average; clustering in the halo would only raise the local density and strengthen the bound). At $T = v_{R}$ the causal minimum exceeds $3.6 \times 10^{-25}$, more than $12$ times the limit with $\pi$ bracketed adversely on both sides (the check pins $100\,(\text{minimum}/\text{maximum})^{2}$ to the integer part $16341$). Since the minimum scales as $T^{3}$, it exceeds the limit for every $T \ge 0.48\,v_{R}$ (checked at $0.48\,v_{R}$), and the check confirms it does not at $0.3\,v_{R}$ with the brackets reversed, so the threshold lies in $(0.3, 0.48)\,v_{R}$; the check pins its sixth power, $(T_{*}/v_{R})^{6}$, to integer parts $6119 \times 10^{-6}$ and $4217 \times 10^{-6}$ with the two bracketings, that is $T_{*} \in (0.40, 0.43)\,v_{R}$ ($0.409\,v_{R} = 1.70 \times 10^{10}$ GeV at the central value of $\pi$). At $T = v_{R}$ it exceeds the limit for every $\beta \ge 10^{-4}$.

(iii) With $M_{M} > 4 \cdot 3 \cdot v_{R}/g_{R} > 10^{12}$ GeV, Preskill's floor is above $10^{-12}$, more than $10^{10}$ times the causal minimum, so annihilation leaves $n_{M}/s \ge \min(\text{initial}, \text{floor})$, which is at least the causal minimum. Closure, $\Omega_{M}h^{2} = M_{M}(n_{M}/s)s_{0}/(\rho_{c}/h^{2}) < 0.12$ with $\rho_{c}/h^{2} = 1.0537 \times 10^{-5}$ GeV cm$^{-3}$, allows $n_{M}/s$ up to $3.9 \times 10^{-22}$ at $M_{M} = 1.1 \times 10^{12}$ GeV, above the minimum, so the flux limit, not closure, is the constraint.

(iv) By Theorem \claimref{synthesis.seeding-remainder-excludes-perturbative-pati-salam}, the seeding origin needs $c\,T_{d}^{4} \ge A/2$; with $c \le 2.28$, which covers the odd towers' vectors ($c = 0.529$) and four towers of Standard-Model content ($c = 2.27$) at the derived production coefficients, $T_{d}^{4} \ge A/(2c)$ gives $T_{d} > 8.3 \times 10^{11}$ GeV (the threshold at $c = 2.28$ lies in $(8.3, 8.4) \times 10^{11}$ GeV), above $20\,v_{R}$. Gauge interactions thermalise the bath at these temperatures, and unless (e)'s exception applies the Pati-Salam symmetry is restored, so the broken phase forms again at $T_{c}$ in the descended phase, a causal process covered by (K): by (ii) the result is excluded unless $T_{c} < 0.48\,v_{R}$ or a Langacker-Pi stage removes the monopoles after they form. No step uses Premise (R) (\claimref{neutrino.r-descent-reheats-below-lower-heavy}) or a bound on the couplings; the stored incompatibility of those two with the seeding origin is therefore not the only obstruction to it. For the seeding temperature to lie below $2.0 \times 10^{10}$ GeV, where (ii) no longer bites, the coefficient would have to exceed $A/(2(2.0 \times 10^{10})^{4}) > 6.8 \times 10^{6}$, more than $2.8 \times 10^{8}$ vectors at the per-vector coefficient $16087\pi^{3}/20736000 = 2.4055 \times 10^{-2}$.

(v) If $T_{d} < T_{c}$, the descent deposits the broken phase directly. If its vacuum orientation is uncorrelated beyond the descent's Hubble length, (K) applies at $T_{d}$ with the same function, so $T_{d} \ge 0.48\,v_{R}$ is excluded. Only a deposit aligned beyond its Hubble length escapes, the same pre-metric correlation that Theorem \claimref{synthesis.finite-range-descent-deposit-cannot-seed-cmb} finds the primordial spectrum would need. Such correlations are available in principle: a deposit correlated by the frozen point's stabiliser depends on the descent slice only through separations relative to $|x_{d}|$, about four present Hubble radii (filed alongside, synthesis.frozen-point-geometry-reads-descent-as-vacuum); its obstruction concerns the spectrum of $\zeta$, not the orientation of the broken vacuum, so it does not close this escape.

\emph{Scope.} The chain's $M_{I}$ and $g_{R}$ are the stored record-level values; a different chain (for example a two-step breaking through $U(1)_{B-L}$, or a much higher $M_{I}$ off the chain) moves the threshold with $v_{R}$. The transition temperature $T_{c}$ is not computed (the store has no Pati-Salam potential); $T_{c}/v_{R} < 0.48$ needs a Higgs quartic small against the gauge coupling. The literature inputs (a) to (e) are from memory and should be imported.

Numerical values (float, not evidence): $s_{0} = 2891.3$ cm$^{-3}$; causal minimum at $v_{R}$, $4.23 \times 10^{-25}$; flux maximum at $\beta = 7 \times 10^{-4}$, $2.90 \times 10^{-26}$; ratio $14.6$; threshold $1.70 \times 10^{10}$ GeV; threshold speed at $T = v_{R}$, $\beta = 4.8 \times 10^{-5}$; $M_{M} \approx 1.1 \times 10^{12}$ GeV.

The check verifies the constant $1/1440$, the hypercharge assignments, tracelessness and period, the phase unit in $\Cl(2,4)$, the bounds on $v_{R}$ and $s_{0}$, the causal minimum against the flux limit with the threshold bracketed in $(0.3, 0.48)\,v_{R}$, the annihilation floor and closure comparison, and the seeding temperature and coefficient.
